\subsection{NODAL VECTORS AND INVERSE ROOTS}

Lemma 2. Let S be a closed subgroup of T and Z(S) its normalizer in G.

\begin{enumerate}
\def\labelenumi{(\roman{enumi})}
\item
  \(\mathrm{R}(\mathrm{Z}(\mathrm{S})_0,\mathrm{T})\) is the set of \(\alpha \in \mathrm{R}(\mathrm{G},\mathrm{T})\) such that \(\alpha (\mathrm{S}) = \{1\}\) .
\item
  The centre of \(\mathrm{Z(S)_0}\) is the intersection of the \(\operatorname {Ker}\alpha\) for \(\alpha \in \mathrm{R}(\mathrm{Z(S)}_0,\mathrm{T})\) .
\item
  If S is connected, Z(S) is connected.
\end{enumerate}

The Lie algebra \(\mathrm{L}(\mathrm{Z}(\mathrm{S}))_{(\mathbf{C})}\) consists of the invariants of S on \(g_{C}\) (Chap. III, §9, no. 3, Prop. 8), and hence is the direct sum of \(t_{C}\) and the \(g^{\alpha}\) for which \(\alpha(\mathrm{S})=\{1\}\) , hence (i). Assertion (ii) follows from Prop. 8 (no. 4), and assertion (iii) has already been proved (§2, no. 2, Cor. 5 of Th. 2).

THEOREM 1. Let \(\alpha \in \mathrm{R}(\mathrm{G},\mathrm{T})\) . The centralizer \(Z_{\alpha}\) of the kernel of \(\alpha\) is a connected closed subgroup of \(\mathrm{G}\) ; its centre is \(\operatorname{Ker} \alpha\) ; its derived group \(\mathrm{D}(Z_{\alpha}) = S_{\alpha}\) is a connected closed semi-simple subgroup of \(\mathrm{G}\) of rank 1. We have \(\mathrm{R}(Z_{\alpha},\mathrm{T}) = \{\alpha, -\alpha\}\) and \(\dim Z_{\alpha} = \dim T + 2\) .

Let \(Z_{\alpha}^{\prime}\) be the centralizer of \((\mathrm{Ker}\alpha)_0\) . By Lemma 2, this is a connected closed subgroup of G, and \(\mathrm{R}(Z_{\alpha}^{\prime},\mathrm{T})\) is the set of \(\beta \in \mathrm{R}(\mathrm{G},\mathrm{T})\) such that \(\beta ((\mathrm{Ker}\alpha)_0) = \{1\}\) . Clearly, \(\{\alpha , - \alpha \} \subset \mathrm{R}(Z_{\alpha}^{\prime},\mathrm{T})\) . Conversely, let \(\beta \in \mathrm{R}(Z_{\alpha}^{\prime},\mathrm{T})\) ; since \((\mathrm{Ker}\alpha)_0\) is of finite index in \(\mathrm{Ker}\alpha\) , there exists an integer \(r\neq 0\) such that \(t^{r\beta} = 1\) for \(t\in \mathrm{Ker}\alpha\) . From the exactness of the sequence

\[
0 \longrightarrow \mathbf {Z} \longrightarrow \mathrm{X(T)} \longrightarrow \mathrm{X(Ker} \alpha) \longrightarrow 0
\]

corresponding by duality to the exact sequence

\[
0 \longrightarrow \operatorname{Ker} \alpha \longrightarrow \mathrm{T} \stackrel {{\alpha}} {{\longrightarrow}} \mathrm{U} \longrightarrow 0,
\]

it follows that \(r\beta\) is a multiple of \(\alpha\) ; by Chap. VIII, §2, no. 2, Th. 2 (i), this implies that \(\beta \in \{\alpha, -\alpha\}\) . Thus, \(\mathrm{R}(Z_{\alpha}', T) = \{\alpha, -\alpha\}\) . It follows (Lemma 2) that the centre of \(Z_{\alpha}'\) is \(\operatorname{Ker} \alpha\) , so \(Z_{\alpha}' = Z_{\alpha}\) . Finally, by Cor. 1 of Prop. 4 (§1, no. 4), \(\mathrm{D}(\mathrm{Z}_{\alpha})\) is a connected closed semi-simple subgroup of \(\mathrm{G}\) ; it is of rank 1 because \(\mathcal{DL}(\mathrm{Z}_{\alpha})_{(\mathbf{C})} = \mathfrak{g}^{\alpha} + \mathfrak{g}^{-\alpha} + [\mathfrak{g}^{\alpha},\mathfrak{g}^{-\alpha}]\) .

COROLLARY. There exists a morphism of Lie groups \(\nu : \mathbf{SU}(2, \mathbf{C}) \to \mathrm{G}\) with the following properties:

\begin{enumerate}
\def\labelenumi{\alph{enumi})}
\item
  The image of \(\nu\) commutes with the kernel of \(\alpha\) .
\item
  For all \(a \in U\) , we have \(\nu\left(\begin{array}{cc}a & 0 \\ 0 & \bar{a}\end{array}\right) \in T\) and \(\alpha \circ \nu\left(\begin{array}{cc}a & 0 \\ 0 & \bar{a}\end{array}\right) = a^{2}\) .
\end{enumerate}

If \(\nu_{1}\) and \(\nu_{2}\) are two morphisms from \(\mathbf{SU}(2,\mathbf{C})\) to \(\mathrm{G}\) with the preceding properties, there exists \(a\in \mathbf{U}\) such that \(\nu_{2} = \nu_{1}\circ \operatorname {Int}\left( \begin{array}{cc}a & 0\\ 0 & \bar{a} \end{array} \right)\) .

By Th. 1 and Prop. 6 of §3, no. 6, there exists a morphism of Lie groups \(\nu : \mathbf{SU}(2, \mathbf{C}) \to \mathrm{S}_{\alpha}\) that is surjective with discrete kernel. Then \(\nu^{-1}(\mathrm{T} \cap \mathrm{S}_{\alpha})\) is a maximal torus of \(\mathbf{SU}(2, \mathbf{C})\) (§2, no. 3, Prop. 1). Since the maximal tori of \(\mathbf{SU}(2, \mathbf{C})\) are conjugate (§2, no. 2, Th. 2), we can assume, replacing \(\nu\) by \(\nu \circ \operatorname{Int}s\) (with \(s \in \mathbf{SU}(2, \mathbf{C})\) ) if necessary, that \(\nu^{-1}(\mathrm{T} \cap \mathrm{S}_{\alpha})\) is the group of diagonal matrices in \(\mathbf{SU}(2, \mathbf{C})\) . Then \(\nu\begin{pmatrix} a & 0 \\ 0 & \bar{a} \end{pmatrix} \in \mathrm{T}\) for all \(a \in \mathbf{U}\) , and the map

\[
\left( \begin{array}{c c} a & 0 \\ 0 & \bar {a} \end{array} \right) \mapsto \alpha \circ \nu \left( \begin{array}{c c} a & 0 \\ 0 & \bar {a} \end{array} \right)
\]

is a root of \(\mathbf{SU}(2,\mathbf{C})\) , and hence is equal to one of the two maps \(\begin{pmatrix}a&0\\ 0&\bar{a}\end{pmatrix}\mapsto a^{2}\) or \(\begin{pmatrix}a&0\\ 0&\bar{a}\end{pmatrix}\mapsto a^{-2}\) (§3, no. 6, formulas (19)). In the first case, the homomorphism \(\nu\) has the required properties; in the second case, the homomorphism \(\nu\circ Int\theta\) has them (loc. cit., formulas (18)).

If \(\nu_{1}\) and \(\nu_{2}\) are two morphisms from \(\mathbf{SU}(2,\mathbf{C})\) to G satisfying the stated conditions, they both map \(\mathbf{SU}(2,\mathbf{C})\) into \(S_{\alpha}\) (condition a)), hence are both universal coverings of \(S_{\alpha}\) . Hence, there exists an automorphism \(\varphi\) of \(\mathbf{SU}(2,\mathbf{C})\) such that \(\nu_{2} = \nu_{1} \circ \varphi\) , and we conclude by using Prop. 9 of no. 4.

It follows from the preceding corollary that the homomorphism \(\nu_{\mathrm{T}}\) from \(\mathbf{U}\) to \(\mathrm{T}\) , defined by \(\nu_{\mathrm{T}}(a) = \nu \left( \begin{array}{cc}a & 0\\ 0 & \bar{a} \end{array} \right)\) for \(a\in \mathbf{U}\) , is independent of the choice of \(\nu\) . Denote by \(K_{\alpha}\in \Gamma (\mathrm{T})\) the image under \(\Gamma (\nu_{\mathrm{T}})\) of the element \(2\pi i\) of \(\Gamma (\mathbf{U}) = 2\pi i\mathbf{Z}\) ; it is called the nodal vector associated to the root \(\alpha\) . We have \(\langle \alpha ,K_{\alpha}\rangle = 2\) , in other words (no. 2, formula (2)) \(\delta (\alpha)(K_{\alpha}) = 4\pi i\) ; since \(K_{\alpha}\) belongs to the intersection of \(\mathfrak{t}\) and the \(\mathrm{L}(\mathrm{S}_{\alpha})_{(\mathbf{C})}\) , we have

\[
K _ {\alpha} = 2 \pi i H _ {\delta (\alpha)},\tag{13}
\]

where \(H_{\delta(\alpha)}\) is the inverse root associated to the root \(\delta(\alpha)\) of \((\mathfrak{g}_{\mathbf{C}}, \mathfrak{t}_{\mathbf{C}})\) (Chap. VIII, §2, no. 2). In other words, when \(\Gamma(\mathrm{T}) \otimes \mathbf{R}\) is identified with the dual of \(\mathrm{X(T)}\otimes \mathbf{R}\) via the pairing \(\langle ,\rangle ,K_{\alpha}\) is identified with the inverse root \(\alpha^{\vee}\in (\mathrm{X(T)}\otimes \mathbf{R})^{*}\) .

Remark. For all \(x \in R\) , we have

\[
\nu \left( \begin{array}{c c} \exp (2 \pi i x) & 0 \\ 0 & \exp (- 2 \pi i x) \end{array} \right) = \nu_ {\mathrm{T}} (e ^ {2 \pi i x}) = \exp (x K _ {\alpha}).\tag{14}
\]

In particular:

\[
\nu \left( \begin{array}{c c} - 1 & 0 \\ 0 & - 1 \end{array} \right) = \nu_ {\mathrm{T}} (- 1) = \exp \Bigl (\frac {1}{2} K _ {\alpha} \Bigr).\tag{15}
\]

It follows that \(\nu\) is injective if and only if \(K_{\alpha} \notin 2\Gamma(T)\) , in other words if there exists \(\lambda \in \mathrm{X}(T)\) such that \(\langle \lambda, K_{\alpha} \rangle \notin 2Z\) . When \(g_{C}\) is simple, \(\nu\) is injective unless \(g_{C}\) is of type \(B_{n}\) , \(C(G) = \{1\}\) and \(\alpha\) is a short root (cf.~Chap. VI, Plates).

In the remainder of this paragraph we denote by \(\mathrm{R}^{\vee}(\mathrm{G},\mathrm{T})\) the set of nodal vectors \(K_{\alpha}\) for \(\alpha\in\mathrm{R}(\mathrm{G},\mathrm{T})\) . This is a subset of \(\varGamma(\mathrm{T})\) that the canonical injection of \(\varGamma(\mathrm{T})\) into \(t_{C}\) identifies with the homothety with ratio \(2\pi i\) of the inverse root system \(\mathrm{R}^{\vee}(\mathfrak{g}_{\mathbf{C}},\mathfrak{t}_{\mathbf{C}})=\{H_{\delta(\alpha)}\}\) of \(\delta(\mathrm{R})\) . It follows that \(\mathrm{R}^{\vee}(\mathrm{G},\mathrm{T})\) generates the R-vector space \(\mathrm{L}(\mathrm{T}\cap\mathrm{D}(\mathrm{G}))\) , and hence that its orthogonal complement in \(\mathrm{X}(\mathrm{T})\) is \(\mathrm{X}(\mathrm{T}/(\mathrm{T}\cap\mathrm{D}(\mathrm{G})))\) .

Denote by \(\mathrm{Aut(T)}\) the group of automorphisms of the Lie group T; the Weyl group \(W = W_{G}(T)\) (§2, no. 5) can be identified with a subgroup of \(\mathrm{Aut(T)}\) . On the other hand, recall (Chap. VIII, §2, no. 2, Remark 4) that the Weyl group \(W(\mathfrak{g}_{\mathbf{C}}, t_{\mathbf{C}})\) of the split reductive algebra ( \(\mathfrak{g}_{\mathbf{C}}, t_{\mathbf{C}}\) ) operates on \(t_{\mathbf{C}}\) , and thus is canonically identified with a subgroup of \(\mathbf{GL}(t_{\mathbf{C}})\) .

PROPOSITION 10. The map \(u \mapsto \mathrm{L}(u)_{(\mathbf{C})}\) from \(\operatorname{Aut}(\mathrm{T})\) to \(\mathbf{GL}(\mathfrak{t}_{\mathbf{C}})\) induces an isomorphism from \(\mathrm{W}\) to the Weyl group of the split reductive Lie algebra \((\mathfrak{g}_{\mathbf{C}}, \mathfrak{t}_{\mathbf{C}})\) . For all \(\alpha \in \mathrm{R}\) , \(\mathrm{W}_{\mathrm{Z}_{\alpha}}(\mathrm{T})\) is of order 2, and the image under the preceding isomorphism of the non-identity element of \(\mathrm{W}_{\mathrm{Z}_{\alpha}}(\mathrm{T})\) is the reflection \(s_{H_{\delta(\alpha)}}\) .

The map under consideration is injective. It remains to show that its image is equal to \(\mathrm{W}(\mathfrak{g}_{\mathbf{C}}, t_{\mathbf{C}})\) .

Let \(g \in \mathrm{N}_{\mathrm{G}}(\mathrm{T})\) . With the notations in Chap. VIII, §5, no. 2, we have \(\operatorname{Ad} g \in \operatorname{Aut}(\mathfrak{g}_{\mathbf{C}}, \mathfrak{t}_{\mathbf{C}}) \cap \operatorname{Int}(\mathfrak{g}_{\mathbf{C}})\) , so \(\operatorname{Ad} g \in \operatorname{Aut}_{0}(\mathfrak{g}_{\mathbf{C}}, \mathfrak{t}_{\mathbf{C}})\) (loc. cit., no. 5, Prop. 11). By loc. cit., no. 2, Prop. 4, the automorphism of \(\mathfrak{t}_{\mathbf{C}}\) induced by \(\operatorname{Ad} g\) belongs to \(\mathrm{W}(\mathfrak{g}_{\mathbf{C}}, \mathfrak{t}_{\mathbf{C}})\) . Thus, the image of \(\mathrm{W}\) in \(\mathbf{GL}(\mathfrak{t}_{\mathbf{C}})\) is contained in \(\mathrm{W}(\mathfrak{g}_{\mathbf{C}}, \mathfrak{t}_{\mathbf{C}})\) .

Let \(\alpha \in \mathrm{R}(\mathrm{G},\mathrm{T})\) , and let \(\nu : \mathbf{SU}(2,\mathbf{C}) \to \mathrm{G}\) be a morphism of Lie groups having the properties in the Cor. of Th. 1. The image under \(\nu\) of the element \(\theta\) of \(\mathbf{SU}(2,\mathbf{C})\) has the following properties (§3, no. 6, formulas (17)):

\[
a) (\operatorname{Int} \nu (\theta)) (t) = t \quad \text { if } t \in \operatorname{Ker} \alpha ,
\]

\[
b) (\operatorname{Int} \nu (\theta)) (t) = t ^ {- 1} \quad \text { if } t \in \mathrm{T} \cap \mathrm{S} _ {\alpha}.
\]

It follows that \(\operatorname{Ad}\nu(\theta)\) induces the identity on \(\operatorname{Ker}\delta(\alpha)\subset\mathfrak{t}_{\mathbf{C}}\) , and induces the map \(x\mapsto-x\) on \([\mathfrak{g}^{\alpha},\mathfrak{g}^{-\alpha}]\) , hence coincides with the reflection \(s_{H_{\delta(\alpha)}}\) . Thus, the image of W contains all the \(s_{H_{\delta(\alpha)}}\) , and hence is equal to \(\mathrm{W}(\mathfrak{g}_{\mathbf{C}}, t_{\mathbf{C}})\) . In particular \(\mathrm{W}_{\mathrm{Z}_{\alpha}}(\mathrm{T})\) is of order 2, and hence consists of the identity and \(\operatorname{Int}\nu(\theta)\) . This completes the proof of the proposition.

COROLLARY. Assume that G is semi-simple. Then every element of G is the commutator of two elements of G.

Let \(c\) be a Coxeter transformation of the Weyl group \(\mathrm{W}(\mathfrak{g}_{\mathbf{C}},\mathfrak{t}_{\mathbf{C}})\) (Chap. V, §6, no. 1), and let \(n\) be an element of \(\mathrm{N}_{\mathrm{G}}(\mathrm{T})\) whose class in \(\mathrm{W}\) is identified with \(c\) by the isomorphism defined in the proposition. Denote by \(f_{c}\) the morphism \(t\mapsto (n,t)\) from \(\mathrm{T}\) to \(\mathrm{T}\) ; for \(x\in \mathfrak{t}_{\mathbf{C}}\) , we have \(\mathrm{L}(f_c)_{(\mathbf{C})}(x) = (\mathrm{Ad}n)(x) - x = c(x) - x\) .

By Th. 1 of Chap. V, §6, no. 2, the endomorphism \(c\) of \(\mathbf{t}_{\mathbf{C}}\) has no eigenvalue equal to 1. Consequently, \(\mathrm{L}(f_c)\) is surjective, and hence so is \(f_c\) . It follows that every element of T is the commutator of two elements of G, which implies the corollary in view of Th. 2, §2, no. 2.

